%-------------------------------------------------------------------------------
%  Recommendation Letter for Xingwei Qu
%  Author : Jie Fu (Fujie) — Research Scientist, IQuest Research (Ubiquant)
%  Style  : Academic / Formal, refined typography
%-------------------------------------------------------------------------------
\documentclass[11pt,a4paper]{letter}

%---------------------------- Packages ----------------------------------------
\usepackage[T1]{fontenc}
\usepackage[utf8]{inputenc}
\usepackage{lmodern}
\usepackage[osf,sc]{mathpazo}              % Palatino w/ old-style figures + small caps
\usepackage[scaled=0.92]{helvet}           % Helvetica for sans accents
\linespread{1.12}
\usepackage{microtype}
\usepackage[a4paper,
            top=1.00in,
            bottom=1.15in,
            left=1.10in,
            right=1.10in]{geometry}
\usepackage{xcolor}
\usepackage{graphicx}
\usepackage{parskip}
\usepackage{enumitem}
\usepackage{ragged2e}
\usepackage{soul}                          % letterspacing

%---------------------------- Colors ------------------------------------------
\definecolor{accent}{HTML}{0B2545}         % Oxford blue
\definecolor{accentSoft}{HTML}{5C7CA8}     % muted blue for hairlines
\definecolor{subtle}{HTML}{4A4A4A}         % warm dark grey for meta text
\definecolor{rule}{HTML}{0B2545}

\usepackage[colorlinks=true,
            urlcolor=accent,
            linkcolor=accent]{hyperref}

%---------------------------- Sender details ----------------------------------
\newcommand{\senderName}{Jie Fu}
\newcommand{\senderTitle}{Research Scientist}
\newcommand{\senderAffiliation}{IQuest Research, Ubiquant}
\newcommand{\senderEmail}{\href{mailto:jie.fu@iquestlab.com}{jie.fu@iquestlab.com}}
\newcommand{\senderWebsite}{\href{https://bigaidream.github.io}{bigaidream.github.io}}
\newcommand{\senderScholar}{\href{https://scholar.google.com/citations?user=66osleIAAAAJ}{Google Scholar}}

%---------------------------- Header macro ------------------------------------
\makeatletter
\renewcommand{\opening}[1]{%
  \begin{flushleft}
    {\fontsize{20}{24}\selectfont\bfseries\color{accent}%
      \so{Jie Fu}\,\normalfont\color{accent}\bfseries, Ph.D.}\\[3pt]
    {\sffamily\small\color{accent}\senderTitle\ \,\textbullet\,\ \senderAffiliation}\\[6pt]
    {\sffamily\footnotesize\color{subtle}%
       \senderEmail \quad\textcolor{accentSoft}{\textbar}\quad
       \senderWebsite \quad\textcolor{accentSoft}{\textbar}\quad
       \senderScholar}
  \end{flushleft}
  \vspace{2pt}
  {\color{rule}\rule{\linewidth}{0.8pt}}\\[-3pt]
  {\color{accentSoft}\rule{\linewidth}{0.3pt}}
  \vspace{18pt}

  \begin{flushleft}
    {\sffamily\small\color{subtle}\@date}
  \end{flushleft}
  \vspace{4pt}

  \noindent #1\par\nobreak
  \vspace{8pt}
}
\makeatother

%---------------------------- Closing macro -----------------------------------
\renewcommand{\closing}[1]{%
  \vspace{14pt}\noindent #1\\[46pt]
  \noindent{\bfseries\color{accent}\senderName}
}

%---------------------------- Document ----------------------------------------
\begin{document}

\setlength{\parindent}{0pt}
\setlength{\parskip}{8pt}

\begin{letter}{}

\date{May 24, 2026}

\opening{Dear Microsoft Team,}

\justifying

I am writing to strongly recommend Xingwei Qu.

I have known Xingwei for many years, and I have had the opportunity to observe
his development over a long period of time. What makes him particularly
impressive is that I have seen him change from a pure engineer into a serious
researcher. When I first knew him, his strongest qualities were execution
speed, coding ability, and engineering ownership. Over the years, especially
through his work in large-scale model training and pretraining research, he
has developed the ability to ask deep research questions and to connect
implementation with scientific insight.

My own background gives me a particular lens on this. I spent roughly five
years in the Mila / University of Montreal research environment, including
postdoctoral work with Yoshua Bengio, and I have been deeply influenced by
research questions around human-friendly, reliable, and safe AI. From that
perspective, I think Xingwei is a very promising candidate. He is not only
interested in building more capable systems; he is increasingly interested in
how these systems reason, how they can be evaluated, and how their internal
computation can be made more efficient, inspectable, and controllable.

Xingwei's engineering ability is genuinely outstanding. He has worked as a
senior algorithm engineer at Xiaopeng and later on large-scale pretraining
and model research at ByteDance Seed. He has the rare ability to take an
ambiguous technical goal, turn it into a concrete implementation plan, build
the infrastructure, debug the system, run experiments, and produce useful
evidence quickly. He is not merely a prototype engineer; he thinks about
training stability, data quality, evaluation, reproducibility, and how the
code supports the research question.

At the same time, his research ability has grown very quickly. His recent
work on Dynamic Large Concept Models, ConceptMoE, latent reasoning, and
compute-efficient pretraining shows that he is now thinking beyond standard
engineering optimization. He is asking questions about the unit of reasoning
in large models, how computation should be allocated, how token-level
representations can be compressed into higher-level semantic units, and how
we can evaluate reasoning beyond surface-level chain-of-thought behavior.
These are exactly the kinds of questions that require both strong systems
intuition and real research taste.

Recently, we have also been discussing and collaborating on AI safety-related
directions, including latent reasoning, verification, agent evaluation, and
how to construct training and evaluation pipelines that can reveal whether
models are following robust internal procedures. I find his motivation
genuine. He is not approaching AI safety as a slogan; he is approaching it as
an engineer-researcher who wants to build concrete methods, benchmarks, and
systems that make advanced AI more reliable and understandable.

The main area where Xingwei can still improve is prioritization. Because he
is extremely energetic and can implement quickly, he sometimes opens too many
research directions at once. I have seen this in recent discussions where he
proposed several promising safety-related directions simultaneously. My
feedback to him was to narrow the work into one sharper claim, one clearer
experimental pipeline, and one concrete set of success and failure criteria.
He responded very well, which is another reason I trust his trajectory. He is
ambitious, but also coachable and willing to refine his thinking.

Overall, I strongly recommend Xingwei. He has the execution velocity of an
excellent engineer, the emerging taste of a strong researcher, and a genuine
interest in the kinds of safety and reliability questions that matter for
advanced AI systems. I believe he would be a very strong fit for Anthropic.

\closing{Sincerely,}

\end{letter}
\end{document}
